\chapter{直线、圆与角}\label{ch:g6:lines}

几何从直尺、三角板和圆规开始。本章
确定基本词汇 --- 点、线段、射线、直线、圆 ---
介绍直线的两种特殊位置（平行与
垂直），并讲授如何度量和画角。

\section{点、线段、直线}

\begin{definition}[基本对象]\label{def:g6:lines:objects}
过两个不同的点 $A$ 和 $B$ 可以画出以下对象：
\begin{itemize}
  \item \emph{线段}\index{线段} $[AB]$：直线上的
        位于 $A$、$B$ 之间的部分（它有长度，记作 $AB$）；
  \item \emph{射线} $[AB)$：从 $A$ 出发，经过 $B$，
        并无限延伸；
  \item \emph{直线}\index{直线} $(AB)$：向两个方向
        无限延伸。两点唯一确定一条直线。
\end{itemize}
同一直线上的点称为\emph{共线}。
\end{definition}

\begin{omfigure}
\begin{tikzpicture}[scale=1.0]
  \fill (0,1.8) circle (1.5pt) node[above, font=\small] {$A$};
  \fill (2.5,1.8) circle (1.5pt) node[above, font=\small] {$B$};
  \draw[thick, omDef] (0,1.8) -- (2.5,1.8);
  \node[right, font=\small] at (2.9,1.8) {线段 $[AB]$};
  \fill (0,0.9) circle (1.5pt) node[above, font=\small] {$A$};
  \fill (2.5,0.9) circle (1.5pt) node[above, font=\small] {$B$};
  \draw[thick, omThm] (0,0.9) -- (4.1,0.9);
  \node[right, font=\small] at (4.3,0.9) {射线 $[AB)$};
  \fill (0,0) circle (1.5pt) node[above, font=\small] {$A$};
  \fill (2.5,0) circle (1.5pt) node[above, font=\small] {$B$};
  \draw[thick, omMeth] (-1.4,0) -- (4.1,0);
  \node[right, font=\small] at (4.3,0) {直线 $(AB)$};
\end{tikzpicture}

{\small 两个相同的点，三种不同的对象。方括号表示
哪一端停止、哪一端延伸：$[$ 表示停止，$($ 表示延伸。}
\end{omfigure}

\section{平行线与垂线}

\begin{definition}[垂直、平行]\label{def:g6:lines:perp}
两条直线相交成直角时，称它们\emph{垂直}\index{垂直}，
记作 $d_1 \perp d_2$，直角处用小方块标出。
若两条直线无论延伸多远都不相交，
则称它们\emph{平行}\index{平行}，
记作 $d_1 \parallel d_2$。
\end{definition}

\begin{omfigure}
\begin{tikzpicture}[scale=1.0]
  \draw[thick] (-1.8,0) -- (1.8,0) node[right, font=\small] {$d_1$};
  \draw[thick, omDef] (0.6,-1.2) -- (0.6,1.4)
    node[above, font=\small] {$d_2$};
  \draw[thin] (0.6,0) ++(-0.25,0) -- ++(0,0.25) -- ++(0.25,0);
  \node[below, font=\small] at (0,-1.5) {$d_1 \perp d_2$};
  \begin{scope}[xshift=6.5cm]
  \draw[thick] (-1.8,-0.5) -- (1.8,0.3) node[right, font=\small] {$d_3$};
  \draw[thick, omThm] (-1.8,-1.3) -- (1.8,-0.5)
    node[right, font=\small] {$d_4$};
  \node[below, font=\small] at (0,-1.5) {$d_3 \parallel d_4$};
  \end{scope}
\end{tikzpicture}

{\small 直角相交（用小方块标出）与两条平行线：
方向相同，且没有交点。}
\end{omfigure}

\begin{proposition}[两个有用的结论]\label{prop:g6:lines:facts}
\begin{enumerate}
  \item 若两条直线都与第三条直线垂直，则它们
        彼此平行。
  \item 若两条直线平行，则与其中一条垂直的每一条直线
        也与另一条垂直。
\end{enumerate}
\end{proposition}
\admitted

\begin{method}[用三角板作图]\label{met:g6:lines:setsquare}
要作直线 $d$ 的垂线，使其过点 $P$：
\begin{enumerate}
  \item 将三角板直角的一条边沿 $d$ 放置；
  \item 沿 $d$ 平移三角板，直至它的另一条边
        到达 $P$；
  \item 沿该边画线，并标出直角。
\end{enumerate}
过 $P$ 作平行线时，先作 $d$ 的垂线，再
过 $P$ 作\emph{该}垂线的垂线
（见\cref{prop:g6:lines:facts}）。
\end{method}

\section{圆}

\begin{definition}[圆]\label{def:g6:lines:circle}
圆心为 $O$、半径为 $r$ 的\emph{圆}\index{圆}是
所有点的集合，这些点的距离恰为 $r$（从 $O$ 量起）。
连接圆心与圆上一点的线段称为\emph{半径}；
连接圆上两点且过圆心的线段称为\emph{直径} ---
其长度为 $2r$；连接圆上两点的线段称为\emph{弦}。
\end{definition}

\begin{omfigure}
\begin{tikzpicture}[scale=1.15]
  \draw[thick] (0,0) circle (1.5);
  \fill (0,0) circle (1.5pt) node[below left, font=\small] {$O$};
  \fill (1.5,0) circle (1.5pt) node[right, font=\small] {$M$};
  \draw[omDef, thick] (0,0) -- (1.5,0)
    node[midway, above, font=\small] {半径};
  \fill (-1.06,1.06) circle (1.5pt) node[above left, font=\small] {$A$};
  \fill (1.06,-1.06) circle (1.5pt) node[below right, font=\small] {$B$};
  \draw[omThm, thick] (-1.06,1.06) -- (1.06,-1.06)
    node[pos=0.25, right=2pt, font=\small] {直径};
  \fill (0.52,1.41) circle (1.5pt) node[above, font=\small] {$C$};
  \fill (1.41,0.52) circle (1.5pt) node[right, font=\small] {$D$};
  \draw[omMeth, thick] (0.52,1.41) -- (1.41,0.52)
    node[midway, above right=-2pt, font=\small] {弦};
\end{tikzpicture}

{\small 一个圆，半径为 $[OM]$，直径为 $[AB]$（过圆心的弦，
长是半径的两倍），还有一条弦 $[CD]$。}
\end{omfigure}

\begin{example}\label{ex:g6:lines:circle}
“画出圆心为 $O$、经过 $A$ 的圆”：把圆规从 $O$ 张开到 $A$，
——半径就是距离 $OA$——然后转动一周。圆上的每一点
到 $O$ 的距离都与 $A$ 相等。
\end{example}

\section{角}

\begin{definition}[角]\label{def:g6:lines:angle}
两条射线 $[AB)$ 和 $[AC)$，端点都是 $A$，构成
\emph{角}\index{角} $\widehat{BAC}$；点 $A$ 是它的
\emph{顶点}（总是中间那个字母！）。角的度量单位是
\emph{度}（$^\circ$），从 $0^\circ$ 到 $360^\circ$ 为一整
周。角可分为：
\begin{itemize}
  \item \emph{直角}，若其度数为 $90^\circ$；
  \item \emph{锐角}，若其度数小于 $90^\circ$；
  \item \emph{钝角}，若其度数介于 $90^\circ$ 与
        $180^\circ$ 之间；
  \item \emph{平角}，若其度数为 $180^\circ$（两条射线构成
        一条直线）。
\end{itemize}
\end{definition}

\begin{omfigure}
\begin{tikzpicture}[scale=0.9]
  \draw[thick] (0,0) -- (1.8,0);
  \draw[thick] (0,0) -- (40:1.8);
  \draw[omDef, ->] (0.7,0) arc (0:40:0.7);
  \node[omDef, font=\small] at (20:1.0) {$40^\circ$};
  \node[below, font=\small] at (0.9,-0.3) {锐角};
  \begin{scope}[xshift=3.6cm]
  \draw[thick] (0,0) -- (1.8,0);
  \draw[thick] (0,0) -- (0,1.8);
  \draw[thin] (0.3,0) -- (0.3,0.3) -- (0,0.3);
  \node[below, font=\small] at (0.9,-0.3) {直角};
  \end{scope}
  \begin{scope}[xshift=7.2cm]
  \draw[thick] (0,0) -- (1.8,0);
  \draw[thick] (0,0) -- (135:1.8);
  \draw[omThm, ->] (0.6,0) arc (0:135:0.6);
  \node[omThm, font=\small] at (65:1.0) {$135^\circ$};
  \node[below, font=\small] at (0.9,-0.3) {钝角};
  \end{scope}
  \begin{scope}[xshift=11.6cm]
  \draw[thick] (-1.5,0) -- (1.5,0);
  \fill (0,0) circle (1.2pt);
  \draw[omMeth, ->] (0.5,0) arc (0:180:0.5);
  \node[below, font=\small] at (0,-0.3) {平角};
  \end{scope}
\end{tikzpicture}

{\small 四类角。直角（$90^\circ$）是
基准：锐角表示更小，钝角表示更大。}
\end{omfigure}

\begin{method}[用量角器测量角]\label{met:g6:lines:protractor}
\begin{enumerate}
  \item 将量角器的中心恰好放在角的
        顶点上；
  \item 使其 $0^\circ$ 刻度线与角的一条边对齐；
  \item 读出角的另一条边所穿过的刻度——
        使用从对齐那条边上的 $0$ 开始的刻度；
  \item 用眼睛做粗略检查：锐角的读数应小于
        $90^\circ$，钝角则应更大。
\end{enumerate}
\end{method}

\begin{example}\label{ex:g6:lines:estimate}
测量之前先估一估！比一张纸的角稍微张开的角，
略大于 $90^\circ$；半个直角是 $45^\circ$；三分之一个直角是
$30^\circ$。如果你的量角器读数为
$150^\circ$，而该角看起来是锐角，说明你读错了刻度，
正确的度数应为 $180^\circ - 150^\circ = 30^\circ$。
\end{example}

\section{练习}

\begin{exercise}[$\star$]\label{exo:g6:lines:1}
画出三个不共线的点 $A$、$B$、$C$。用不同颜色画出：
线段 $[AB]$、射线 $[CA)$、直线 $(BC)$。
\end{exercise}

\begin{exercise}[$\star$]\label{exo:g6:lines:2}
判断下列说法是否正确：“$[AB]$ 与 $[BA]$ 是同一条线段。”
“$[AB)$ 与 $[BA)$ 是同一条射线。”“$(AB)$ 与 $(BA)$ 是
同一条直线。”每个答案都要说明理由。
\end{exercise}

\begin{exercise}[$\star$]\label{exo:g6:lines:3}
画一条直线 $d$ 和一点 $P$（不在 $d$ 上）。用三角板
作图：作 $d$ 的垂线（过点 $P$），再作 $d$ 的平行线，
同样过 $P$。描述你的作图步骤。
\end{exercise}

\begin{exercise}[$\star$]\label{exo:g6:lines:4}
直线 $a$ 与 $b$ 都垂直于直线 $c$，又有第四条直线
$e$ 垂直于 $a$。关于 $a$ 与 $b$，你能得出什么结论？
关于 $e$ 与 $c$ 呢？请用 \cref{prop:g6:lines:facts} 说明理由。
\end{exercise}

\begin{exercise}[$\star$]\label{exo:g6:lines:5}
画一个圆心为 $O$、半径为 $3$ 厘米的圆。在圆上取一点 $M$，
在圆内取一点 $N$，在圆外取一点 $P$。
关于距离 $OM$、$ON$、$OP$ 与 $3$ 厘米的大小关系，你能得出什么结论？
\end{exercise}

\begin{exercise}[$\star$]\label{exo:g6:lines:6}
一个圆的直径为 $9$ 厘米，它的半径是多少？另一个圆的半径为
$2.6$ 厘米，它的直径是多少？
\end{exercise}

\begin{exercise}[$\star$]\label{exo:g6:lines:7}
用三个字母记出图中标出的角，再对它分类（锐角、直角、
钝角、平角）：$72^\circ$ 的角，顶点为 $R$，两边分别射向 $S$ 和 $T$；
$148^\circ$ 的角，顶点为 $B$，两边分别射向 $A$ 和 $C$；
$90^\circ$ 的角，顶点为 $O$，两边分别
射向 $M$ 和 $N$。
\end{exercise}

\begin{exercise}[$\star$]\label{exo:g6:lines:8}
先估计，再用量角器量出你自己画的三角形的
三个角的度数。把三个角的度数相加，你发现了什么？
（请把这一结论留给\cref{ch:g7:triangles}使用。）
\end{exercise}

\begin{exercise}[$\star$]\label{exo:g6:lines:9}
用量角器画一个角 $\widehat{xOy}$，度数为 $65^\circ$，再画一个
度数为 $130^\circ$ 的角。不用量角器，怎样由第一个角
得到第二个角？
\end{exercise}

\begin{exercise}[$\star\star$]\label{exo:g6:lines:10}
地图上画有两个村庄 $A$ 和 $B$。求满足下列条件的点在哪里：
距离 $3$ 厘米（指到 $A$ 的距离）\emph{且}距离 $2$ 厘米（指到 $B$ 的距离）？请画图说明
这样的点可能有多少个（$0$、$1$ 或 $2$，取决于
距离 $AB$）。
\end{exercise}

\begin{exercise}[$\star\star$]\label{exo:g6:lines:11}
钟面指向 3 点时，两根指针之间的夹角是多少？对于 5 点和
6 点时，同样的问题又是怎样的？（一整圈
是 $360^\circ$，对应 $12$ 小时。）
\end{exercise}

\begin{exercise}[$\star\star\star$]\label{exo:g6:lines:12}
画一条线段 $[AB]$，长 $6$ 厘米。仅用圆规作出点 $C$，使得
$AC = BC = 6$ 厘米，并量出角 $\widehat{CAB}$ 的大小。
你作出的是什么三角形？关于它的各个角，
你有什么猜想？
\end{exercise}

\section{问题：钟面上的几何}

\begin{problem}[{周末问题 --- 用整圆的几分之几表示角：在钟面上读角，
寻找两根指针相遇的时刻，
以及把一天切成一张饼图}]
\label{pb:g6:lines:1}
钟面就是一个能报时的量角器：它的表盘是一整圈
$360^\circ$，被十二个整点刻度分成十二个
相等的角。\Cref{exo:g6:lines:11}中已测量了指针在
3 点和 6 点时的夹角；本题将建立完整的理论，
--- 包括\emph{两根指针都不}指向刻度的时刻，
并回答一个很少有成年人答对的问题（“两根指针
多久会恰好重叠在一起？”），最后把一整天
切成一张饼图。

\medskip
\noindent\textbf{第一部分 --- 整圆的几分之几。}
\begin{enumerate}
  \item 一整圈是 $360^\circ$。半圈、四分之一圈、
        十二分之一圈分别是多少度？
  \item 十二个整点刻度把钟面分成圆心处十二个
        相等的角。相邻两个刻度之间是多少度？
        用数刻度的方法，再次求出
        \cref{exo:g6:lines:11}中 3 点时的答案。
  \item 求 1 点、
        4 点和 7 点时两根指针之间的夹角（7 点时，两根指针
        把钟面分成两个角；“两根指针之间的角”
        总是指\emph{较小}的那一个。）
  \item 分针在 $60$ 分钟内转一整圈。分针每分钟
        扫过多少度？
  \item 时针从一个刻度走到下一个刻度 ---
        $30^\circ$ --- 用时 $60$ 分钟。那么
        \emph{它}每分钟扫过多少度？（这是一个小数，
        见\cref{ch:g6:decimals}。）
\end{enumerate}

\medskip
\noindent\textbf{第二部分 --- 没有指针指向
刻度的时刻。}
\begin{enumerate}[resume]
  \item 3:30 时，分针指向 $6$。时针究竟
        在哪里？分别算出两根指针与 $12$ 的夹角
        （按顺时针方向量），并由此推出两根指针
        之间的夹角。
  \item 6:30 时，很多人猜两根指针会重叠在一起。
        先猜一猜，再像第 6 题那样算出夹角。
        谁猜对了？
  \item 计算 9:15 时两根指针之间的夹角。
  \item 请说明：在 1:00 到 1:10 之间，
        两根指针一定恰好重叠。1:00 时分针相对时针
        在哪里？1:10 时又在哪里？
        哪根指针走得更快？为什么这一点
        足以说明问题？
  \item 在十二小时内，两根指针恰好重叠
        在一起的次数是多少？（相邻两个整点之间
        各有一次相遇 --- 但有一个例外：
        11 点到 12 点之间会发生什么？请列出
        近似的相遇时刻并计数。）
        一整天内又有多少次？
\end{enumerate}

\medskip
\noindent\textbf{第三部分 --- 围绕一点的角：
饼图。}
\begin{enumerate}[resume]
  \item 钟面上的十二个扇形合起来充满整个钟面，
        它们的角之和为 $360^\circ$。请说明为什么
        对\emph{任意}一组角都成立，这些角有公共顶点，
        绕该顶点恰好填满一整圈，
        既不重叠，也没有空隙。
  \item 佐伊记录了她 $24$ 小时的一天：睡觉 $9$ 小时，上学
        $6$ 小时，玩耍 $3$ 小时，用餐 $2$ 小时，其他活动 $4$ 小时。
        她想画一个\emph{饼图}：一个圆盘，每项活动占一个扇形，
        扇形的角与时间成正比例。
        每项活动占一天的几分之几？其扇形的圆心角是多少度？
        检验这五个角的和是否等于
        $360^\circ$。
  \item 请逐步说明，如何用圆规和量角器
        画出佐伊的饼图
        （见\cref{met:g6:lines:protractor}）：第一条半径画在何处，
        以及每个新扇形如何从前一个扇形
        的终止处开始。
  \item 另一张描述一天的饼图中，一个扇形恰好是
        $90^\circ$。它占整个圆盘的几分之几，
        它代表多少小时？
  \item 压轴题，回到钟面：在 $24$ 小时内，时针
        转了多少整圈？分针呢？秒针呢？
        （其中一个答案超过一千。）
\end{enumerate}
\end{problem}
